\begin{thebibliography}{25}
\bibitem[1]{Arikiq} Susumu Ariki, \emph{Graded q-Schur algebras}, 2009, \url{https://arxiv.org/abs/0903.3453}.
\bibitem[2]{AK} Susumu Ariki and Kazuhiko Koike, \emph{A Hecke algebra of $(\mathbb{Z}/r\mathbb{Z})\wr S_n$ and construction of its irreducible representations}, Adv.\ Math.\ 106 (1994), no.~2, 216--243.
\bibitem[3]{BKKL} Jonathan Brundan and Alexander Kleshchev, \emph{Blocks of cyclotomic Hecke algebras and Khovanov-Lauda algebras}, Invent.\ Math.\ 178 (2009), no.~3, 451--484.
\bibitem[4]{DJM} Richard Dipper, Gordon James, and Andrew Mathas, \emph{Cyclotomic q-Schur algebras}, Math.\ Z.\ 229 (1998), no.~3, 385--416.
\bibitem[5]{DMmorita} Richard Dipper and Andrew Mathas, \emph{Morita equivalences of Ariki--Koike algebras}, Math.\ Z.\ 240 (2002), no.~3, 579--610.
\bibitem[6]{GTL} Sachin Gautam and Valerio Toledano Laredo, \emph{Yangians and quantum loop algebras}, Selecta Math.\ (N.S.) 19 (2013), no.~2, 271--336.
\bibitem[7]{GGOR} Victor Ginzburg, Nicolas Guay, Eric Opdam, and Rapha\"el Rouquier, \emph{On the category $\mathcal{O}$ for rational Cherednik algebras}, Invent.\ Math.\ 154 (2003), no.~3, 617--651.
\bibitem[8]{GreenSchur} Richard M.~Green, \emph{The affine q-Schur algebra}, J.\ Algebra 215 (1999), no.~2, 379--411.
\bibitem[9]{HMsemi} Jun Hu and Andrew Mathas, \emph{Seminormal forms and cyclotomic quiver Hecke algebras of type A}, Math.\ Ann.\ 364 (2016), no.~3-4, 1189--1254.
\bibitem[10]{Jimbo} Michio Jimbo, \emph{A q-analogue of $U(\mathfrak{gl}(N+1))$, Hecke algebra, and the Yang-Baxter equation}, Lett.\ Math.\ Phys.\ 11 (1986), no.~3, 247--252.
\bibitem[11]{KLI} Mikhail Khovanov and Aaron D.~Lauda, \emph{A diagrammatic approach to categorification of quantum groups. I}, Represent.\ Theory 13 (2009), 309--347.
\bibitem[12]{KLII} Mikhail Khovanov and Aaron D.~Lauda, \emph{A diagrammatic approach to categorification of quantum groups II}, Trans.\ Amer.\ Math.\ Soc.\ 363 (2011), no.~5, 2685--2700.
\bibitem[13]{Lusgraded} George Lusztig, \emph{Affine Hecke algebras and their graded version}, J.\ Amer.\ Math.\ Soc.\ 2 (1989), no.~3, 599--635.
\bibitem[14]{MacBourbaki} Ian G.~Macdonald, \emph{Affine Hecke algebras and orthogonal polynomials}, in S\'eminaire Bourbaki : volume 1994/95, expos\'es 790-804, Ast\'erisque, no.~237, Soci\'et\'e math\'ematique de France, 1996, Exp.\ No.~797, pp.~189--207.
\bibitem[15]{MacHecke} Ian G.~Macdonald, \emph{Affine Hecke algebras and orthogonal polynomials}, Cambridge Tracts in Mathematics, vol.~157, Cambridge University Press, Cambridge, 2003.
\bibitem[16]{MS18} Ruslan Maksimau and Catharina Stroppel, \emph{Higher level affine Schur and Hecke algebras}, \url{https://arxiv.org/abs/1805.02425}.
\bibitem[17]{MiSt} Vanessa Miemietz and Catharina Stroppel, \emph{Affine quiver Schur algebras and p-adic $GL_n$}, Selecta Math.\ (N.S.) 25 (2019), no.~32.
\bibitem[18]{Rou2KM} Rapha\"el Rouquier, \emph{2-Kac-Moody algebras}, \url{https://arxiv.org/abs/0812.5023}.
\bibitem[19]{SWschur} Catharina Stroppel and Ben Webster, \emph{Quiver Schur algebras and q-Fock space}, \url{https://arxiv.org/abs/1110.1115}.
\bibitem[20]{Uglov} Denis Uglov, \emph{Canonical bases of higher-level q-deformed Fock spaces and Kazhdan-Lusztig polynomials}, in Physical combinatorics (Kyoto, 1999), Progr.\ Math., vol.~191, Birkh\"auser, 2000, pp.~249--299.
\bibitem[21]{Webalt} Ben Webster, \emph{Representation theory of the cyclotomic Cherednik algebra via the Dunkl-Opdam subalgebra}, \url{https://arxiv.org/abs/1609.05494}.
\bibitem[22]{WebwKLR} Ben Webster, \emph{Weighted Khovanov-Lauda-Rouquier algebras}, to appear in Documenta Mathematica, \url{https://arxiv.org/abs/1209.2463}.
\bibitem[23]{Webmerged} Ben Webster, \emph{Knot invariants and higher representation theory}, vol.~250, Mem.\ Amer.\ Math.\ Soc., no.~1191, American Mathematical Society, 2017.
\bibitem[24]{WebRou} Ben Webster, \emph{Rouquier's conjecture and diagrammatic algebra}, Forum Math.\ Sigma 5 (2017), e27 (71 pages).
\bibitem[25]{Kbook} Charles A.~Weibel, \emph{The K-book. An introduction to algebraic K-theory}, Graduate Studies in Mathematics, vol.~145, American Mathematical Society, Providence, RI, 2013.
\end{thebibliography}
